\section{Experiments}
\label{sec:experiments}

\subsection{Setup}

Unless otherwise noted, all experiments use \textbf{Qwen2.5-1.5B-Instruct}
\cite{qwen2.5} as the base LLM, extracting hidden states at layer 28 (the
final transformer block output). The nine stories per class come from the
pre-computed cache at
\texttt{line4/emotion\_vectors/qwen\_1.5B/stories\_*.json}.
For bootstrap confidence intervals, we use 1{,}000 resamples.
All datasets are loaded via the \texttt{datasets} library at the versions
current in April~2026.

\subsection{Zero-Shot Sentiment Classification}

\paragraph{SST-2.}
We project the 872 SST-2 validation sentences onto the V-axis and threshold at
the median (yielding a zero-label binary classifier). We also report a
calibrated threshold tuned on a held-out 200-sample subset.

\begin{table}[h]
\centering
\caption{SST-2 zero-shot sentiment classification. 95\% bootstrap CIs in
brackets. Our method matches a supervised logistic regression trained on
80--100 labeled SST-2 samples.}
\label{tab:sst2}
\begin{tabular}{lcc}
\toprule
Method & AUC & Accuracy \\
\midrule
Random direction (mean of 5) & 0.509 & 0.509 \\
SBERT MiniLM-L6 + same pipeline & 0.713 & 0.658 \\
Supervised LR, $N=10$ labels & 0.680 $\pm$ 0.058 & 0.608 $\pm$ 0.059 \\
Supervised LR, $N=100$ labels (L28 features) & 0.877 & 0.778 \\
Supervised LR, $N=5000$ labels (L28 features) & 0.923 & 0.848 \\
Supervised LR, $N=67{,}349$ (full SST-2 train, L28) & 0.951 & 0.878 \\
Direct LLM prompting (greedy) & -- & 0.928 \\
\rowcolor{blond}
\textbf{\ours (zero labels)} & \textbf{\sstauc{}} [0.844, 0.890] & \textbf{\sstacc{}} \\
\rowcolor{blond}
\textbf{\ours (200-sample threshold calib.)} & \sstauc{} & 79.17\% \\
\bottomrule
\end{tabular}
\end{table}

Our method at AUC $0.868$ sits between the supervised classifier with 50
labeled examples and the one with 100. It uses zero labels.

\paragraph{Four More Sentiment Benchmarks.}
We apply the same V-axis to IMDB (balanced 500+500), Yelp Polarity
(balanced 500+500), Rotten Tomatoes (1{,}066), and SemEval-2017 Twitter
sentiment (balanced 500+500).

\begin{table}[h]
\centering
\caption{Same V-axis applied to four additional sentiment benchmarks. All
zero-shot (no labels).}
\label{tab:benchmarks}
\begin{tabular}{lccc}
\toprule
Benchmark & $n$ & AUC & Median-thresh Acc \\
\midrule
IMDB (balanced) & 1{,}000 & \imdbauc{} & \imdbacc{} \\
\textbf{Yelp Polarity (balanced)} & 1{,}000 & \textbf{\yelpauc{}} & \textbf{\yelpacc{}} \\
Rotten Tomatoes & 1{,}066 & \mrauc{} & 76.55\% \\
Twitter SemEval-17 & 1{,}000 & \twauc{} & 83.80\% \\
\bottomrule
\end{tabular}
\end{table}

All five benchmarks use the \textit{identical} nine-story-centroid V-axis.
The same direction generalizes across review styles (Yelp restaurants, IMDB
movies, SST movies), lengths (SST-2 sentences vs.\ IMDB paragraphs), and
registers (Twitter informal vs.\ financial formal text).

\subsection{Word-Level Sentiment Lexicons}

\paragraph{Hu \& Liu 6{,}789 words.}
We extract the residual-stream embedding of each word in a minimal template
\texttt{"The word \textbackslash"\{w\}\textbackslash" means"}, then project onto
the V-axis.

\begin{table}[h]
\centering
\caption{Word-valence lexicon recovery. Pearson $|r|$ between our
V-axis projection and human-annotated valence scores.}
\label{tab:lexicon}
\begin{tabular}{lccc}
\toprule
Lexicon & $n$ & Pearson $|r|$ & Random baseline \\
\midrule
Hu \& Liu (binary +1/$-$1) & 6{,}789 & \textbf{\huliur{}} [0.74, 0.78] & 0.12 \\
Warriner \textit{et al.\@} (continuous 1--9) & 3{,}000 & \textbf{\warrinerr{}} & 0.07 \\
NRC 32-word fallback (curated) & 32 & 0.95 & 0.09 \\
\bottomrule
\end{tabular}
\end{table}

On both lexicons the signal-to-noise ratio vs.\ random projection is
$\geq6\times$. On Hu \& Liu the 95\% bootstrap confidence interval excludes
zero tightly ($[0.74, 0.78]$).

\paragraph{Continuous Sentence-Level Valence (EmoBank).}
EmoBank \cite{buechel2017emobank} provides 10{,}062 English sentences with
continuous VAD scores. Our V-axis projection correlates at
$|r|=\emobankr{}$ [0.485, 0.512] with ground-truth valence scores.
A supervised Ridge regression trained on the full feature space reaches
$|r|=0.66$; our method recovers $72\%$ of that signal with no labels.

\subsection{Fine-Grained Sentiment (SST-5)}

SST-5 provides 5-class ordinal valence (very negative \ldots very positive).
Our V-axis projection correlates at Pearson $r = 0.620$ (Spearman $0.627$)
with the 0--4 label on the 2{,}210-sample test set, confirming that the
axis captures continuous valence not just sign.

\subsection{Data Efficiency: How Few Stories Are Enough?}

\begin{figure}[h]
\centering
\includegraphics[width=0.7\textwidth]{nshot_curve}
\caption{Data-efficiency curve. \emph{Top}: SST-2 AUC as a function of the
number of stories per class (our method, zero labels). \emph{Bottom}:
supervised LR trained on $N$ labeled SST-2 samples. Our zero-label method
matches the supervised curve at $N=80$--$100$.}
\label{fig:nshot}
\end{figure}

A single story per class (9 total samples) already gives AUC 0.74 on SST-2.
Fifteen stories per class (135 total) match the supervised 5{,}000-label LR.

\begin{table}[h]
\centering
\caption{N-shot ablation on SST-2. Mean over 3 subsampling seeds.}
\label{tab:nshot}
\begin{tabular}{cc|cc}
\toprule
\multicolumn{2}{c}{\ours (zero labels)} & \multicolumn{2}{c}{Supervised LR ($N$ labels)} \\
$n$ per class & SST-2 AUC & $N$ labels & SST-2 AUC \\
\midrule
1 & 0.740 $\pm$ 0.043 & 10 & 0.680 $\pm$ 0.058 \\
3 & 0.738 $\pm$ 0.014 & 50 & 0.834 $\pm$ 0.035 \\
9 & 0.781 $\pm$ 0.029 & 100 & 0.878 $\pm$ 0.012 \\
15 & \textbf{0.831 $\pm$ 0.010} & 500 & 0.906 $\pm$ 0.003 \\
25 & 0.833 $\pm$ 0.008 & 2000 & 0.919 $\pm$ 0.001 \\
50 & 0.837 $\pm$ 0.000 & 5000 & 0.914 $\pm$ 0.000 \\
\bottomrule
\end{tabular}
\end{table}

At $n=1$ story per class (9 total LLM-generated samples with no labels),
we already beat the supervised classifier at $N=10$ labeled samples.

\subsection{The V-Shape: Layer-wise Emergence of the Valence Direction}

\begin{figure}[h]
\centering
\includegraphics[width=0.9\textwidth]{vshape_qwen15B}
\caption{Per-layer SST-2 AUC of the V-axis extracted at each layer,
for Qwen2.5-1.5B (28 layers). Valence is present at layer~1 (embedding-level
signal), absent through layers 2--26 (the dilution zone), and
re-crystallizes at layers 27--28.}
\label{fig:vshape}
\end{figure}

\begin{table}[h]
\centering
\caption{Layer-wise SST-2 AUC for Qwen2.5-1.5B. The valence direction
re-crystallizes abruptly at layers 27--28.}
\label{tab:layerwise}
\begin{tabular}{lccccc}
\toprule
Layer & 1 & 14 & 26 & 27 & 28 \\
\midrule
SST-2 AUC & \textbf{0.752} & 0.525 & 0.532 & \textbf{0.836} & \textbf{0.837} \\
EmoBank V $|r|$ & 0.370 & 0.005 & 0.002 & 0.517 & 0.496 \\
\bottomrule
\end{tabular}
\end{table}

Cross-layer V-axis alignment (cosine similarity of the V-axes at different
layers): layer 1 vs.\ 28: 0.29; layer 14 vs.\ 28: \textbf{0.07} (nearly
orthogonal); layer 27 vs.\ 28: 0.95. The final-layer valence direction is
\emph{not} inherited from earlier layers---it is newly constructed by the
last 1--2 blocks.

\subsection{Cross-Family Universality}

\begin{table}[h]
\centering
\caption{Same method, five LLM families. EmoBank V $|r|$ at the final layer.}
\label{tab:crossmodel}
\begin{tabular}{lccc}
\toprule
Model & Family & $|V_L|$ & Recovery of Full+Ridge baseline \\
\midrule
Qwen2.5-0.5B-Instruct & Qwen & 0.485 & -- \\
\textbf{Qwen2.5-1.5B-Instruct} & Qwen & \textbf{0.475} & 72\% \\
Qwen2.5-7B-Instruct & Qwen & 0.455 & \textbf{77\%} \\
\midrule
Pythia-1.4b & GPT-NeoX & 0.461 & -- \\
Bloom-560m & BLOOM & 0.342 & -- \\
\bottomrule
\end{tabular}
\end{table}

All five models produce a usable late-layer valence direction. The
magnitude varies (Qwen strongest, BLOOM weakest) but every model shows the
same V-shape pattern.

\paragraph{Functional alignment across models.}
Different models have different hidden-state dimensionalities, so we
cannot directly compare their V-axis \emph{vectors}. Instead, we
project the same 81 stories onto each model's V-axis, producing a
per-story valence score for every model. After a simple sign alignment
(positive emotions on the positive side), we compute pairwise cosine
similarity and Spearman rank correlation of these 81-dimensional
per-story valence vectors.

\begin{table}[h]
\centering
\caption{Functional cross-model agreement on per-story valence
projections. Despite different architectures and sizes, all four models
rank the 81 stories nearly identically on valence.}
\label{tab:crossmodel_func}
\begin{tabular}{lcc}
\toprule
Model pair & Cosine sim & Spearman $\rho$ \\
\midrule
Qwen2.5-0.5B $\times$ Qwen2.5-1.5B & \textbf{0.980} & \textbf{0.977} \\
Qwen2.5-0.5B $\times$ Pythia-1.4b & 0.974 & 0.972 \\
Qwen2.5-0.5B $\times$ Bloom-560m & 0.944 & 0.935 \\
Qwen2.5-1.5B $\times$ Pythia-1.4b & \textbf{0.984} & \textbf{0.982} \\
Qwen2.5-1.5B $\times$ Bloom-560m & 0.944 & 0.932 \\
Pythia-1.4b $\times$ Bloom-560m & 0.937 & 0.924 \\
\midrule
\textbf{Mean (all 6 pairs)} & \textbf{0.961} & \textbf{0.954} \\
\bottomrule
\end{tabular}
\end{table}

Cross-model mean pairwise cosine is $0.961$ and Spearman $\rho$ is
$0.954$, with a minimum of $0.937$ (Pythia-1.4b vs.\ Bloom-560m). The
four models --- two Qwen sizes, a GPT-NeoX model, and a BLOOM model ---
functionally agree on which of the 81 stories are more or less positive,
even though they have different architectures, training data, and
internal coordinate systems. This is stronger evidence for
universality than the per-model $|V_L|$ numbers in
\cref{tab:crossmodel}, which only show that each model individually has
a usable V-axis. The high cosine across language-model families
suggests that late-layer valence direction is a convergent solution of
the next-token prediction objective rather than an artifact of any
single training pipeline.

\begin{figure}[h]
\centering
\includegraphics[width=0.9\textwidth]{vshape_crossmodel}
\caption{V-shape pattern on three model families. Top: Qwen2.5-1.5B (28
layers), middle: Pythia-1.4b (24 layers), bottom: Bloom-560m (24 layers).
All three show an early-layer signal (L1--L4), a middle-layer dilution
zone, and a final-layer re-crystallization.}
\label{fig:crossmodel}
\end{figure}

\subsection{Cross-Language Generalization}

A central claim of this paper is that the valence axis transfers to
non-English text. We evaluate two transfer protocols. \emph{Protocol 1
(transfer).} Extract the V-axis from the 81 \emph{English} stories at
Qwen2.5-1.5B L28. Score held-out 64-sentence story splits in English,
Spanish, French, German, and Chinese (no re-extraction, no re-fitting, no
translation). \emph{Protocol 2 (within-language cosines).} Extract a
separate V-axis per language (15 stories per emotion) and measure pairwise
cosines. The first protocol is the deployment-relevant question; the
second tests whether the underlying direction is language-invariant at
the representation level.

\begin{figure}[h]
\centering
\includegraphics[width=0.85\textwidth]{fig_crosslang}
\caption{English V-axis transfers zero-shot to four non-English languages.
Bars show binary-sentiment AUC and median-threshold accuracy on held-out
64-sentence story splits; the English axis is reused with no translation
and no re-extraction. Chance line at $0.5$.}
\label{fig:crosslang}
\end{figure}

\begin{table}[h]
\centering
\caption{Cross-lingual transfer of the English V-axis (Qwen2.5-1.5B L28,
zero labels). The same axis classifies non-English held-out stories
well above chance without any multilingual fine-tuning or translation.}
\label{tab:crosslang_transfer}
\begin{tabular}{lccc}
\toprule
Target language & $n$ & Transfer AUC & Transfer Acc. \\
\midrule
English (in-domain) & 64 & 0.990 & 93.75\% \\
Spanish & 64 & 0.949 & 90.63\% \\
French & 64 & 0.991 & 96.88\% \\
German & 64 & 0.972 & 90.63\% \\
\rowcolor{blond}
\textbf{Chinese} & 64 & \textbf{1.000} & \textbf{100.00\%} \\
\bottomrule
\end{tabular}
\end{table}

The English axis reaches AUC $\geq 0.949$ on all four non-English test
splits, including Chinese at $1.000$. Cross-lingual transfer does not
require a multilingual encoder or translation; the multilingual LLM's
late-layer residual stream appears to place valence along a direction
that is shared across the languages it was trained on.

\begin{table}[h]
\centering
\caption{Within-language V-axis cosine similarity (Protocol 2).
Valence axes share roughly half their direction across languages; arousal
axes are essentially orthogonal.}
\label{tab:crosslang}
\begin{tabular}{lcc}
\toprule
Pair & V-axis cosine & A-axis cosine \\
\midrule
en vs.\ es & \textbf{0.466} & 0.186 \\
es vs.\ fr & \textbf{0.462} & $-$0.249 \\
de vs.\ en & 0.310 & 0.239 \\
de vs.\ es & 0.273 & 0.114 \\
en vs.\ fr & 0.162 & $-$0.097 \\
de vs.\ fr & 0.229 & 0.023 \\
\bottomrule
\end{tabular}
\end{table}

Within-language cosines are positive but modest ($[0.16, 0.47]$), indicating
that each language's 9-centroid geometry finds a direction partially but
not identically aligned with English's. The stronger transfer result
(Table~\ref{tab:crosslang_transfer}) is what matters for deployment: the
English axis is already a good cross-lingual probe in practice, even if
the underlying per-language directions differ by roughly $60^\circ$.

\subsection{Principal-Angle Analysis of the 9-Centroid Geometry}
\label{sec:principal_angle}

We verify that PC1 of the nine centroids is not a trivial
high-variance-axis artifact. The singular values of the mean-centered
centroid matrix at L28 are $(29.13, 15.21, 12.11, 10.17, 9.02, 6.61,
6.27, 4.96, \approx 0)$; PC1 captures $55.8\%$ of the centroid variance
and the ratio of the first two singular values is $1.91$, well separated.
The principal angle between PC1 and PC2 is $90^\circ$ (orthogonal by
construction); the principal angle between PC1 and a uniform random
direction in $\mathbb{R}^{1536}$ is $88.81^\circ \pm 0.83^\circ$ over
$100$ samples. The centroid PC1 is therefore a geometrically principled
direction: it is the unique single-direction summary that maximises
between-category spread, and it is not close to any random direction in
the high-dimensional hidden space.

\subsection{Model-Scale Stability}
\label{sec:scale_stability}

We extend the cross-family evaluation (Table~\ref{tab:crossmodel}) with
a scale sweep on the Qwen2.5 family (0.5B, 1.5B, 7B) at matched
extraction settings. All three scales produce a V-axis with EmoBank
V~$|r| \in [0.485, 0.494]$. The 1.5B and 7B models give essentially
identical valence correlation ($|r| = 0.494$ on both), indicating that
the recipe is not sensitive to model scale in the range
most commonly deployed.

\subsection{Valence/Arousal Asymmetry}
\label{sec:va_asymmetry}

\begin{table}[h]
\centering
\caption{Valence vs.\ arousal extraction success. Same method (PCA of 9 centroids)
applied to both. $|r|$ against EmoBank's continuous V and A scores.}
\label{tab:va_asymmetry}
\begin{tabular}{lcc}
\toprule
Layer & Valence $|r|$ & Arousal $|r|$ \\
\midrule
L1 & 0.370 & 0.012 \\
L14 & 0.005 & 0.013 \\
L27 & \textbf{0.517} & 0.007 \\
L28 & \textbf{0.496} & 0.044 \\
\bottomrule
\end{tabular}
\end{table}

Across every layer, every model family, and every continuous-regression
test, the valence direction is extractable and the arousal direction is not.
The maximum arousal $|r|$ we observed at any layer in any model is 0.16
(Warriner lexicon at L28), still well below the valence number.

\begin{figure}[h]
\centering
\includegraphics[width=0.85\textwidth]{va_asymmetry}
\caption{Valence vs.\ arousal layer-wise on Qwen2.5-1.5B. Valence (green)
has early-layer signal at L1, a middle dilution zone, and a sharp
re-crystallization at L27--L28 reaching $|r|=0.50$. Arousal (red) never
exceeds $|r|=0.05$ at any layer --- the asymmetry is structural, not a
small-effect question.}
\label{fig:va_asym}
\end{figure}

\subsection{PCA-of-Centroids vs.\ Mean-Difference Contrastive Axis}
\label{sec:contrastive}

The standard representation-engineering recipe
\cite{zou2023representation, rimsky2024steering, arditi2024refusal}
extracts a concept direction as the mean-difference between activations
on a pair of contrastive prompts. For a fair comparison we re-partition
our own stories into a positive pool (\texttt{amusement, inspiration, joy,
tenderness}, 80 stories) and a negative pool (\texttt{anger, fear, disgust,
sadness}, 80 stories), compute the mean-difference axis at L28, and
evaluate on the same 872 SST-2 validation sentences.

\begin{table}[h]
\centering
\caption{PCA-of-nine-centroids vs.\ mean-difference contrastive axis.
Both use the same L28 hidden states; the only difference is how the
axis is computed from the story embeddings.}
\label{tab:contrastive}
\begin{tabular}{lc}
\toprule
Axis extraction method & SST-2 val AUC \\
\midrule
Mean-diff contrastive (pos\,$-$\,neg, 80+80 stories) & 0.753 \\
\rowcolor{blond}
\textbf{PCA of 9 centroids (ours)} & \textbf{0.868} \\
\midrule
\emph{Absolute improvement} & \textbf{$+$0.115 AUC} \\
\bottomrule
\end{tabular}
\end{table}

The two axes have absolute cosine $0.527$ --- substantially aligned but
not identical. PCA of nine class-centroids improves SST-2 AUC by $0.115$
over the contrastive-pair baseline on matched features, while also
\emph{removing} the requirement that the operator specify a binary
positive/negative split. In the multilingual setting (Table
\ref{tab:crosslang_transfer}) the PCA axis also transfers without any
analogous contrastive-pair specification across languages.

\subsection{Counterfactual: Is the V-Axis the Only Direction?}
\label{sec:counterfactual}

To quantify precisely what the V-axis contributes, we perform an
activation-ablation study in the spirit of
\cite{arditi2024refusal, park2024linear}. At each of L24, L27, L28 we
(i)~compute the 5-fold cross-validated AUC of a full-rank $L_2$
logistic regression on the L28 hidden features, (ii)~project the V-axis
out of the features (orthogonal complement), and (iii)~refit the same
full-rank LogReg on the ablated features.

\begin{figure}[h]
\centering
\includegraphics[width=0.85\textwidth]{fig_counterfactual}
\caption{V-axis ablation at L24/L27/L28 on SST-2. The blue bar is the
median-thresh 1-D V-axis probe; green is full-rank supervised LogReg
on the hidden state; orange is the same LogReg after projecting the
V-axis out. The change from green to orange is at most $0.003$ AUC,
indicating that valence is distributed across many directions and the
V-axis we extract is a \emph{summary}, not a bottleneck.}
\label{fig:counterfactual}
\end{figure}

\begin{table}[h]
\centering
\caption{Counterfactual V-axis ablation on SST-2 (5-fold CV AUC,
Qwen2.5-1.5B). The full-rank LogReg baseline is the ceiling reachable by
a supervised linear probe on the full hidden state; ``ablated'' is the
same LogReg after the 1-D V-axis is projected out.}
\label{tab:counterfactual}
\begin{tabular}{lccc}
\toprule
Layer & 1-D V-axis (ours) & Full-rank LogReg & After V-axis ablation \\
\midrule
L24 & 0.675 & 0.916 & 0.914 \\
L27 & 0.857 & 0.911 & 0.908 \\
L28 & \textbf{0.868} & 0.911 & 0.910 \\
\bottomrule
\end{tabular}
\end{table}

At L28 the full-rank supervised probe reaches AUC $0.911$; removing our
1-D V-axis from the same features changes that number by $0.001$ ($0.910$).
Valence is therefore richly encoded across many directions of the residual
stream, not localised on a single bottleneck axis. Our contribution is
explicitly framed accordingly: \ours{} produces the most useful
one-dimensional summary of valence extractable under a zero-label
constraint. A supervised baseline that has access to labels and uses all
$1\,536$ dimensions is stronger ($+0.04$ AUC); the axis is a compressed
interpretability-oriented probe, and the counterfactual makes its scope
honest and explicit.

\subsection{PCA Beats Supervised LDA for OOD Transfer}

We compared two ways to extract the concept direction from the story
embeddings: (i)~PCA of the 9 centroids (our method, unsupervised);
(ii)~supervised linear discriminant analysis (LDA) on the per-story vectors
(450 samples, 9 classes).

\begin{table}[h]
\centering
\caption{PCA beats LDA for OOD valence transfer.}
\label{tab:pca_vs_lda}
\begin{tabular}{lcc}
\toprule
Method & EmoBank V $|r|$ \\
\midrule
PCA (unsupervised, 9 centroids) & \textbf{0.498} \\
LDA (supervised, 450 per-story vectors) & 0.119 \\
\bottomrule
\end{tabular}
\end{table}

LDA overfits the 450 training stories and produces a direction that does not
transfer OOD. PCA, despite being unsupervised, is more robust.

\subsection{Speed Benchmark}

\begin{table}[h]
\centering
\caption{Probing vs.\ prompting inference speed (per SST-2 sample,
Qwen2.5-1.5B on V100).}
\label{tab:speed}
\begin{tabular}{lc}
\toprule
Method & ms/sample \\
\midrule
\ours (batched forward pass + projection) & 16.07 \\
Direct prompting (serial, greedy decode) & 123.58 \\
\midrule
\rowcolor{blond}
\textbf{Speedup} & \textbf{\speedup{}} \\
\bottomrule
\end{tabular}
\end{table}

Probing is \speedup{} faster than prompting. The advantage grows with batch
size (probing batches naturally; autoregressive prompting is serial).

\subsection{Baseline Comparison on Standard Sentiment Tasks}

\begin{table}[h]
\centering
\caption{Comprehensive baseline comparison. Our method sits between direct
prompting and SBERT/random baselines.}
\label{tab:allbaselines}
\begin{tabular}{lccc}
\toprule
Method & SST-2 AUC & IMDB AUC & Hu \& Liu $|r|$ \\
\midrule
Random projection (mean of 10) & 0.509 & 0.508 & 0.12 \\
SBERT (MiniLM-L6, same pipeline) & 0.713 & 0.722 & 0.52 \\
\rowcolor{blond}
\textbf{\ours (Qwen2.5-1.5B, zero labels)} & \textbf{0.868} & \textbf{0.895} & \textbf{0.76} \\
Supervised LR ($N=5000$) & 0.914 & -- & -- \\
Direct prompting (greedy) & 0.928 & 0.860 & 0.840 \\
\bottomrule
\end{tabular}
\end{table}

\subsection{Scope of the Method}
\label{sec:scope}

The experiments above define the regime in which \ours{} is useful.
Three complementary bounds --- each quantified, not just asserted --- round
out this picture.

\paragraph{Concept class: global-polarity concepts.}
Repeating the 9-centroid recipe with 15 toxic and 15 safe stories produces
a toxicity direction that reaches AUC $0.59$ on Jigsaw Civil Comments
(near chance). Unlike valence, toxicity depends on specific lexical
triggers whose presence cannot be summarised by a single global direction
that was itself learned from \emph{non-toxic} narrative stories. The
method works when the concept has a clean global-polarity linear
manifestation; we quantify where this breaks (see Supplementary).

\paragraph{Arousal as a scope statement.}
Applying the same PCA recipe to the second principal component does not
recover a usable arousal axis: across every layer and every model we
tested, EmoBank arousal correlation stays at $|r|\leq 0.16$. Combined with
Section~\ref{sec:va_asymmetry} this bounds the operating regime of our
method to \emph{valence-like global concepts that align with the linear
representation hypothesis} \cite{park2024linear}; arousal either lives on a
different, non-linear geometry in the residual stream or is encoded
non-locally. This defines the axis of applicability precisely rather than
treating it as a failure.

\paragraph{Supervised and prompt baselines.}
Direct few-shot prompting reaches $92.78\%$ accuracy on SST-2,
$15$ percentage points above our median-threshold probe, while a
supervised full-rank logistic regression on the same hidden states reaches
AUC $0.911$ ($+0.043$ vs.\ ours). The counterfactual ablation
(Section~\ref{sec:counterfactual}) shows that our V-axis is not the only
direction carrying valence --- removing it barely changes a full-rank
supervised probe. \ours{} therefore sits below these baselines on raw
accuracy \emph{by construction}, because it is a compressed zero-label
1-D probe rather than a full-capacity classifier. The value proposition
is in efficiency, multilingual transfer, interpretability, and continuous
scores rather than in accuracy SOTA.

\paragraph{Cross-LLM ensembling.}
A sign-aligned mean projection across Qwen-0.5B, Qwen-1.5B, Pythia-1.4B,
and Bloom-560m underperforms the best single LLM (Qwen-1.5B) by
$0.02$--$0.04$ AUC on SST-2/IMDB/Yelp. Weaker extractors pull the
ensemble average toward a less-aligned direction; the practical
recommendation is to pick the best single extractor rather than ensemble.
We make the axes themselves available so downstream users can choose.
